\section{Strategic Dynamics: The Theoretical Anchor}
\label{sec:dynamics}

\subsection{From static to dynamic: repeated game dynamics}
\label{subsec:repeated}

A static (one-shot) game ends after a single round of choices. A
repeated game embeds the same stage game in time, so the same producers
choose outputs in every period and discount future profits at a common
factor $\delta \in (0,1)$. The discount factor captures how much present
weight players place on future payoffs: a high $\delta$ corresponds to
patient producers, a low $\delta$ to impatient ones. The reason this
matters is that repetition introduces a credible mechanism of reward and
punishment, which transforms the set of sustainable outcomes relative to
the static game \citep{friedman1971,fudenbergmaskin1986}. Before we
evaluate the cooperative equilibrium, we first study the simplest
possible dynamic adjustment rule: myopic best-response with inertia.
``Myopic'' means each producer optimises against last period's actions
only, without forward-looking calculation; ``inertia''
$\lambda \in (0,1)$ means producers adjust gradually rather than jumping
to the new best response, which is a realistic constraint for an
industry in which drilling, leasing, and contracting take time.

Formally, each producer updates its quantity according to
\begin{equation}
q_i(t+1) \;=\; (1 - \lambda)\, q_i(t) + \lambda\, \mathrm{BR}_i\bigl(q_{-i}(t)\bigr),
\label{eq:3}
\end{equation}
where $\mathrm{BR}_i(\cdot)$ is the static Cournot best response of
player~$i$ to the output vector $q_{-i}(t)$ of all other players.
Starting from an out-of-equilibrium initial condition with $P(0) = 108$,
the price converges monotonically to the static Nash level $P = 60$
within approximately eight periods at $\lambda = 0.2$; the standard
deviation of the simulated price path is 8.14 (see Figure~19). The
dynamics are governed by the dominant eigenvalue $\mu = 1 - 2\lambda$ of
the linearized system around equilibrium, which is why convergence is
monotone for $\lambda \in (0, 0.5)$, critical at $\lambda^* = 0.5$, and
oscillatory for $\lambda \in (0.5, 1)$ (see Figure~20). The economic
content is that even purely myopic adjustment, with no forward-looking
calculation, takes the market to the static Nash equilibrium, which then
becomes the relevant benchmark against which cooperative deviations and
punishments are measured.

Conditional on a large negative shock, defined here as any quarter in
which the simulated price falls more than \$21/bbl below the cartel
price of \$80, the mean recovery path traces out an exponential
approach back to the cooperative level, with a half-life of 5.2 quarters
and asymptotic mean of \$78.7, statistically indistinguishable from the
cartel target (see Figure~21). The 80 percent confidence band around
this mean is wide in the first three periods (where idiosyncratic shock
magnitudes still dominate) and narrows after period 10 once the
punishment phase has cleared. The empirical analogue is the 2020 OPEC+
collapse, where the front-month price recovered approximately 80 percent
of its pre-shock level within five months ($\approx 1.5$ quarters in
our discretisation), slightly faster than our simulated half-life,
consistent with the unusually aggressive April 2020 production cut that
we do not model explicitly. The figure therefore quantifies the speed
at which a credible cooperative arrangement can absorb a large adverse
shock, which is the dynamic counterpart of the static Folk-Theorem
result that we develop in the next subsection.

\subsection{Cartel, punishment, and the Folk Theorem}
\label{subsec:folk}

We now allow producers to consider cooperative quantities that improve
on the Nash outcome. A cartel is a coalition of producers that
coordinates output to maximise joint profit, just as a hypothetical
multi-plant monopolist would. The cartel quota allocation we implement
maximises joint profit subject to proportional output cuts; the price
rises from \$60 to \$80 and every player's per-period payoff rises
relative to Nash. The economic question is not whether this outcome is
attractive (it is, ex ante) but whether it is self-enforcing once
players are free to deviate.

The classical tool for answering this question is the Folk Theorem
\citep{friedman1971,fudenbergmaskin1986}. The Folk Theorem states that,
in an infinitely repeated game with sufficiently patient players, any
individually rational and feasible payoff vector can be supported as a
subgame-perfect equilibrium of the repeated game. The mechanism is that
the threat of future punishment deters present deviation, provided the
discount factor is high enough. The simplest implementation is the
grim-trigger strategy: every player cooperates as long as everyone else
cooperated in every previous period; the first observed deviation
triggers a permanent reversion to the static Nash equilibrium. The
incentive-compatibility (IC) constraint requires that, in every period,
the discounted value of continuing to cooperate exceeds the one-period
gain from deviating plus the discounted value of the post-deviation
punishment:
\begin{equation}
\frac{\pi_i^{\text{coop}}}{1 - \delta} \;\ge\; \pi_i^{\text{dev}} + \frac{\delta\, \pi_i^{\text{nash}}}{1 - \delta},
\label{eq:4}
\end{equation}
which rearranges to
\begin{equation}
\delta \;\ge\; \delta_i^* \;=\; \frac{\pi_i^{\text{dev}} - \pi_i^{\text{coop}}}{\pi_i^{\text{dev}} - \pi_i^{\text{nash}}}.
\label{eq:5}
\end{equation}
The quantity $\delta_i^*$ is the critical discount factor for player~$i$:
the smallest level of patience that makes cooperation incentive
compatible. Our simulations yield $\delta_{\text{US}}^* = 0.455$,
$\delta_{\text{OPEC}}^* = 0.529$, and $\delta_{\text{RUS}}^* = 0.441$
(see Figure~22). OPEC carries the binding constraint because its cost
advantage gives it the highest deviation temptation; the cartel is
sustainable as long as the common discount factor exceeds 0.529. At our
calibrated $\delta = 0.95$, cooperation is comfortably supportable.

We illustrate the mechanism with a simulated deviation by OPEC in period
$t = 5$, followed by ten periods of Nash punishment and a return to
cooperation. OPEC still earns 5,368.67 more under ``cooperate, deviate
once, be punished, recover'' than under ``permanent Nash play'' (see
Figure~23). The implication is that grim-trigger does more than make
cooperation feasible; it makes cooperation strictly the dominant
long-run strategy even when deviations occur. This result connects
directly to the empirical compliance literature on OPEC.
\citet{burbidge1997} show that chronic overproduction by members under
fiscal stress is consistent with the temptation structure that
equation~(\ref{eq:5}) formalizes. \citet{almutairi2025} in \appref{app:almutairi}
provide counterfactual simulations indicating that oil price volatility
would roughly double without OPEC+ discipline, and \citet{montant2025}
finds that OPEC+ exerted stronger price influence in 2017 to 2022 than
in earlier periods but that the post-2022 environment weakened
compliance. The 2020 OPEC+ collapse and its rapid recovery within five
weeks is a textbook example of an off-path punishment phase followed by
re-coordination, and persistent overproduction by Iraq and Kazakhstan
throughout 2024 to 2026 is consistent with players whose fiscal
break-even raises their individual $\delta_i^*$ close to the operative
discount factor. \citet{abreu1986} further demonstrates that brief,
severe stick-and-carrot punishments dominate permanent reversion, a
refinement we exploit only implicitly through the symmetry of the
grim-trigger path.

\subsection{Correlated equilibrium (CE)}
\label{subsec:ce}

Nash and Folk Theorem analysis treats every player as choosing
independently. In many real institutions, however, there is a trusted
mediator who recommends actions confidentially. A correlated equilibrium
(CE) generalises Nash equilibrium to this setting \citep{aumann1974}: a
mediator draws an action profile from a joint distribution on the action
space and privately tells each player its prescribed action; the
distribution is a correlated equilibrium if every player finds it
incentive-compatible to follow the recommendation, conditional on the
information conveyed. The set of correlated equilibria contains all Nash
equilibria and is generally strictly larger; this is the technical
reason a benevolent mediator can do strictly better than uncoordinated
play. The OPEC Secretariat is a natural empirical analogue of the
Aumann mediator, since it issues production targets that members can
either follow or deviate from.

We discretise the action space (ten quantity levels per player) and
solve a linear program over the joint distribution on action profiles,
under three objective functions: maximum total welfare, maximum joint
profit, and maximum-min individual profit. The results sit between the
static Nash and the full cartel: the max-welfare CE delivers
$P = \$60.00$ and $W = 5{,}766.67$ (a gain of about 2 percent over
Nash); the max-joint-profit CE produces $P = \$61.67$ and
$W = 5{,}694.44$; the max-min CE delivers $P = \$60.00$ and
$W = 5{,}600.00$, with the lowest-earning player guaranteed a floor
(see Figures~24, 25 and~26). The CE never reaches the cartel price of
\$80, which reflects the IC constraints that pure cartel quotas
violate.

Two cautions are warranted. The first concerns the distribution of the
welfare gain across players. An allocation is Pareto-improving relative
to a benchmark if no player is worse off and at least one is strictly
better off; this is the standard fairness criterion for changes that
require voluntary agreement. The max-welfare CE fails this test.
Although total welfare rises by about 2 percent above Nash, the gain is
achieved by reallocating production toward the lowest-cost producer:
OPEC's per-period profit rises by 17 percent ($1{,}600 \to 1{,}867$),
while US profit falls by 11 percent ($225 \to 200$) and Russia's by 20
percent ($625 \to 500$). Under an unanimity rule, the relevant decision
rule for a voluntary coalition with no legal enforcement, the
max-welfare CE would therefore be rejected by both US and Russia. Only
the max-min CE, which explicitly imposes a profit floor for the
worst-off player, is Pareto-improving relative to Nash; its welfare
gain is correspondingly smaller ($W = 5{,}600$ versus $5{,}766.67$ for
max-welfare) because protecting the weakest player constrains how
aggressively the mediator can push the equilibrium toward efficiency.
The empirical UAE exit in 2026, Angola in 2024, and Qatar in 2019
(analysed formally through the Shapley-core lens in the next section)
all map onto violations of exactly this Pareto-Nash floor: when
allocated quotas push individual profits below standalone values, the
affected members rationally leave the coalition. The second caution is
narrower and numerical: the discretisation is coarse (ten quantity
levels per player), which produces small support sizes (one to two
profiles in most cases) and unusually clean numerical answers that
should be read as indicative rather than exact.

\subsection{Coalition formation and Shapley values}
\label{subsec:shapley}

Cooperative game theory, in contrast to the non-cooperative theory used
above, asks how the joint surplus generated by cooperation should be
divided when binding agreements are possible. The central object is the
characteristic function $v(S)$, defined for every subset $S$ of the
player set $N$. The number $v(S)$ is the maximum payoff that coalition
$S$ can guarantee itself, irrespective of what non-members do. We
compute $v(S)$ by letting members of $S$ coordinate their outputs to
maximise joint profit while non-members play their simultaneous Nash
best-response. The grand coalition $N$ delivers $v(N)$, the largest
cooperative surplus.

Three concepts complete the framework. The Shapley value
\citep{shapley1953} is the unique allocation that satisfies four axioms
(efficiency, symmetry, null-player, and additivity). Each player's
Shapley payoff $\phi_i$ is its average marginal contribution
$v(S \cup \{i\}) - v(S)$ across all possible orderings of $N$. The core
is the set of allocations such that no sub-coalition can credibly
secede; an allocation is in the core when, for every coalition
$S \subseteq N$, $\sum_{i \in S} x_i \ge v(S)$. Core stability is
therefore the cooperative-game-theoretic analogue of the IC constraint
in repeated games. The broader framework is developed in
\citet{shenoy1979}.

Our characteristic function yields $v(\emptyset) = 0$,
$v(\{\text{US}\}) = 300$, $v(\{\text{OPEC}\}) = 2{,}133.33$,
$v(\{\text{RUS}\}) = 833.33$,
$v(\{\text{OPEC, US}\}) = 2{,}278.12$,
$v(\{\text{RUS, US}\}) = 1{,}012.50$,
$v(\{\text{OPEC, RUS}\}) = 2{,}628.12$, and $v(N) = 3{,}600.00$. The
Shapley allocation is $\phi_{\text{US}} = 477.95$,
$\phi_{\text{OPEC}} = 2{,}202.43$, and $\phi_{\text{RUS}} = 919.62$,
summing to $v(N)$ (see Figure~27). We verify core stability: no
sub-coalition can credibly threaten to break away because, for every
coalition $S \subseteq N$, the sum of Shapley payoffs to the members of
$S$ exceeds the coalition value $v(S)$.

The framework provides a clean theoretical lens on coalition dynamics
in OPEC+. The classical core-stability condition
$\phi_i \ge v(\{i\})$ is equivalent to the statement that each member's
Shapley allocation under prevailing quotas dominates its standalone
alternative. The UAE's withdrawal from OPEC, notified on 28~April 2026 and effective
1~May 2026 \citep{emirates2026}, maps directly onto a violation of this
inequality: Abu Dhabi's capacity of 4.85 mbd was held roughly
30 percent below potential by OPEC+ quotas \citep{woodmac2026}, implying
that the quota-driven allocation undershot the standalone value of
unconstrained capacity. Angola in 2024 and Qatar in 2019 follow the
same logic. \citet{okulloreynes2016} estimate an OPEC cooperation
parameter between Cournot-Nash and perfect cartel, and
\citet{escrihuela2018} construct a coefficient of cooperation for
asymmetric duopoly; both are consistent with the partial-but-stable
allocation we recover here. \citet{ghoddusi2017} provides a useful
caveat: capacity itself is an enforcement mechanism, because members
without spare capacity cannot deviate, and compliance is therefore
partly structural rather than purely strategic.

\subsection{Stochastic demand and Green and Porter monitoring}
\label{subsec:greenporter}

So far, we have assumed that producers perfectly observe each other's
output. In reality, the market price is observable instantly, but
individual production volumes arrive with lag and are revised
repeatedly (BP, OPEC, and IEA estimates routinely diverge by hundreds
of thousands of barrels per day for the same month). This information
asymmetry creates a monitoring problem: a low observed price could
reflect either a genuine deviation by a rival or simply a negative
demand shock that no producer caused. \citet{greenporter1984}
formalized this in their celebrated paper on noncooperative collusion
under imperfect price information. In their model, producers sustain
cooperation through a trigger-price strategy: whenever the observed
price falls below a pre-set threshold $\bar{p}$, all producers revert to
Nash punishment for a fixed number of periods, regardless of whether the
drop was caused by a defection. The equilibrium therefore features
periodic price wars triggered by adverse demand realizations, not by
actual cheating. \citet{porter1983} provided the first empirical
validation of this mechanism using data from the Joint Executive
Committee railroad cartel, and \citet{almoguera2011} in \appref{app:almoguera}
estimated regime-switching probabilities between collusive and Cournot
phases for OPEC over 1974 to 2004.

We model demand uncertainty in two ways. The first is a standard AR(1)
shock to the intercept $a$ with persistence $\rho = 0.6$ and standard
deviation $\sigma = 8$ (in dollars per barrel), which captures the
business-cycle component of demand fluctuations. The second is a
\citet{merton1976} jump-diffusion specification that overlays a compound
Poisson process on the AR(1) diffusion. Jump-diffusion was originally
introduced in financial economics to capture sudden discontinuities in
asset prices; we use it here because oil markets are also subject to
abrupt discontinuities such as pipeline explosions, sanctions
escalations, and armed conflicts. We simulate 300 Monte Carlo paths
over 200 quarters and recover a mean terminal price of \$59.77 with a
P10 to P90 spread of 6.27 around the static Nash benchmark. Scenario
analysis varying $\Delta a$ from $-30$ to $+30$ confirms the linearity
of the static comparative statics: at $\Delta a = -30$ (a severe demand
contraction), OPEC profit falls from 1,600 to 1,056.25 (see Figure~28).

Implementing the Green-Porter trigger with $\bar{p} = 68$ (the midpoint
between the Nash and cartel prices), we recover the regime-switching
pattern that \citet{greenporter1984} predict theoretically and that
\citet{porter1983} and \citet{almoguera2011} document empirically (see
Figures~29 and~30). Table~1 summarises the steady-state statistics
under both demand specifications. Two qualitative readings emerge.
First, under smooth AR(1) demand the cartel is in its cooperative
regime 94 percent of the time and punishment episodes are rare; the
modelled OPEC+ operates roughly as it does in calmer periods such as
2017 to 2019. Second, once we add discontinuous jumps to capture event
risk: pipeline disruptions, sanctions escalations, armed conflicts,
the cooperation fraction falls by nine percentage points, the average
cooperative spell shortens by a third (from 15 to 10 years), and the
number of false punishment triggers more than doubles (291 to 685 over
the same simulation horizon). The punishment-phase length, by contrast,
is essentially invariant across specifications because it is set
exogenously by the trigger rule rather than by the shock process. The
economic content is that the type of uncertainty matters as much as
its magnitude: a market that absorbs smooth volatility can still fail
to absorb a single large jump, which is exactly the pattern observed
in March 2020 (oil-price war coincident with COVID demand collapse)
and in February 2026 (Hormuz closure). \citet{barskykilian2004} argue
that oil price volatility is driven as much by demand-side shocks as
by supply coordination failures, which is precisely the channel through
which false punishment phases arise in our simulation.

We then re-derive the Folk Theorem under stochastic demand. The
deterministic critical discount factor is $\delta^* = 0.53$. Under
uncertainty, with shocks expanding the deviation temptation in good
states and contracting punishment severity in bad ones, it rises to
$\delta^* = 0.60$. Both lie comfortably below our calibrated
$\delta = 0.95$, so cooperation survives (see Figures~31 to~34). We
interpret the 0.07 increment as the ``price of uncertainty'' in cartel
maintenance: a quantitative measure of how much additional patience is
required to keep cooperation robust to noise. \citet{app1990} and
\citet{flm1994} prove the limiting efficiency result that motivates
this exercise, and \citet{jin1999} shows that even private aggregate
information suffices, which is the empirically relevant case for the
oil market because country-level production data arrive late and
revised, but the global spot price is observable in real time.

\subsection{Evolutionary game theory}
\label{subsec:evo}

The final framework of this section treats the producer population as
a biological-style ecosystem of strategy types whose shares evolve over
time. Evolutionary game theory, introduced by \citet{maynardsmith1973}
and developed by \citet[p.~??]{weibull1995}, \citet[p.~??]{hofbauersigmund1998}, and
\citet[p.~??]{hofbauerweibull1996},
differs from classical game theory in two ways. First, it does not
require players to be rational; strategies are inherited or imitated,
and successful strategies spread because their carriers earn higher
payoffs. Second, it predicts long-run population composition rather
than one-shot best responses. The central solution concept is the
evolutionarily stable strategy (ESS): a strategy that, if adopted by
the entire population, cannot be invaded by a small share of mutants
playing any alternative. The dynamic that delivers ESS is the
replicator equation, which states that the share of each strategy grows
in proportion to the gap between its own payoff and the population
average. We use both the two-strategy and three-strategy versions of
the replicator dynamic below.

The two-strategy game lets each player choose Cooperate (C) or Defect
(D). The payoffs satisfy the Prisoner's Dilemma ordering
$\phi^{\text{dev}} = 2{,}025$, $\phi^{\text{coop}} = 1{,}800$,
$\phi^{\text{nash}} = 1{,}600$, and $\phi^{\text{sucker}} = 1{,}500$,
where the sucker payoff is what a cooperator earns against a defector.
Under replicator dynamics the population converges to the all-Defect
corner; the all-Cooperate corner is unstable. Without an enforcement
mechanism, evolutionary selection therefore reproduces the classical
result that cooperation collapses.

We then enrich the strategy set with a third type, Punish (P), which
cooperates on path but retaliates against defectors. The phase diagram
of the three-strategy replicator (see Figures~35, 36 and~37) is
bistable: if the initial share of punishers is sufficiently large, the
population converges to a stable C-P mixture in which cooperation is
sustained; otherwise, defection dominates. The economic intuition is
that a critical mass of credible enforcers is required for cooperation
to survive selection, and the size of that critical mass depends on the
punishment payoff structure.

Sensitivity analysis on the severity of punishment yields a clear
pattern. Moderate punishment multipliers (1.05 to 1.8 times the Nash
payoff) sustain cooperation; excessive punishment (above approximately
1.9 times Nash) destroys the punishers themselves, because retaliation
becomes costlier than the gain from discipline; unconditional,
non-targeted punishment never sustains cooperation, with defection
dominating at every multiplier (see Figure~38). The mapping to OPEC+ is
direct: a moderate, credible threat of flooding (analogous to Saudi
Arabia's 2020 production surge and to the post-2014 episode that
\citet{beharritz2017} analyse in \appref{app:behar-ritz}) supports cooperation,
while indiscriminate or excessive flooding erodes the political capital
required to sustain coordination. \citet{wood2016} provide an
agent-based analogue of this result for the OPEC-versus-Seven-Sisters
transition, and \citet{xinzhang2023} extend the replicator framework to
a tripartite Russia, sanctioning country, and non-sanctioning country
game, with strategies converging to stable states under sanctions.

This section therefore closes the loop on the static-environment
results of the previous one. Repeated interaction, stochastic
monitoring, coalitional bargaining, and evolutionary selection all
converge on the same qualitative finding: cooperation is fragile but
supportable, the binding constraint is OPEC's deviation temptation,
and the threshold conditions have direct empirical counterparts in
observable OPEC+ behaviour. The remainder of the thesis builds on this
anchor by introducing model-free learners (Section~\ref{sec:marl}), an empirical
validation against the 1985, 2014, and 2020 price wars, and a synthesis
of policy implications.
